\documentclass[11pt,a4paper]{article}
\usepackage[margin=1in]{geometry}
\usepackage{amsmath,amssymb,amsfonts,amsthm}
\usepackage{booktabs}
\usepackage{array}
\usepackage{hyperref}
\usepackage{xcolor}

\hypersetup{colorlinks=true,linkcolor=blue,citecolor=blue,urlcolor=blue}

\newcommand{\CP}{\mathbb{CP}}
\newcommand{\Tr}{\mathrm{Tr}}

\newtheorem{theorem}{Theorem}
\theoremstyle{definition}
\newtheorem{definition}[theorem]{Definition}
\theoremstyle{remark}
\newtheorem{remark}[theorem]{Remark}

\begin{document}

\title{Charged Fermion Masses from $\CP^2\times S^3$ Geometry:\\
Nine Predictions, One Input, 1.4\% Mean Accuracy}

\author{Gary Thomas Alcock}

\date{March 2026}

\maketitle

%==========================================================================
% ABSTRACT
%==========================================================================

\begin{abstract}
Within the Density Field Dynamics (DFD) framework, all nine charged
fermion masses follow from a single formula
$m_f = A_f\,\alpha^{n_f}\,v/\!\sqrt{2}$,
where the fine-structure constant $\alpha$ is derived from the
Chern--Simons level sum on $\CP^2$ and the exponents $n_f\in\{0,1,\tfrac{3}{2},\tfrac{5}{2}\}$
arise from spin$^c$ bundle degrees.
The dimensionless prefactors $A_f$ are rational or algebraic numbers
determined by explicit Yukawa operators constructed from $A_5$ conjugacy
class geometry, with no fitted coefficients.
All nine mass \emph{ratios} are therefore parameter-free predictions
of the $\CP^2\times S^3$ topology.

For absolute masses, we present two formulations:
\textbf{(i)}~using the Fermi constant $G_F$ as a single
empirical input ($v/\!\sqrt{2} = 174.1$~GeV), all nine masses are
reproduced with 1.42\% mean error and one free parameter; or
\textbf{(ii)}~using the derived relation $v = M_P\alpha^8\sqrt{2\pi}$,
which introduces a 0.05\% systematic shift in $v$ but eliminates
the last free parameter, yielding 1.40\% mean error with zero
adjustable parameters.
In either formulation, every predicted mass lies within 3.4\%
of its measured value.
We present both formulations explicitly, assess the derivation
status of each ingredient, and identify the open problems whose
resolution would elevate the mass predictions to the same
theorem-grade footing as the $\alpha$ derivation.
\end{abstract}

%==========================================================================
\section{Introduction}
\label{sec:intro}
%==========================================================================

\subsection{The mass problem in perspective}

The Standard Model contains nine charged fermion masses spanning five
orders of magnitude, from $m_e = 0.511$~MeV to $m_t = 172.76$~GeV.
These masses are determined by nine independent Yukawa couplings that
are free parameters of the theory---experimentally measured but not
derived.
Explaining the fermion mass hierarchy remains one of the central
open problems in particle physics.

Within the DFD framework, the fine-structure constant
$\alpha^{-1} = 137.036$ has been derived from first principles:
a single closed-form expression involving the spin$^c$ index on
$\CP^2$, with zero free parameters and agreement to all measured
digits~\cite{DFD_alpha}.
This result establishes that the $\CP^2\times S^3$ topology of the
DFD internal manifold encodes at least some of the Standard Model's
coupling structure.
The natural question is: does the same geometry determine the fermion
masses?

\subsection{What this paper claims---and what it does not}

We derive all nine charged fermion masses from a master formula
\begin{equation}\label{eq:master}
  m_f = A_f \times \alpha^{n_f} \times \frac{v}{\sqrt{2}},
\end{equation}
where each ingredient has a specific geometric origin within
the $\CP^2\times S^3$ framework.
Before proceeding, we state the epistemic status of each component
with full transparency:

\begin{center}
\renewcommand{\arraystretch}{1.25}
\begin{tabular}{lll}
\toprule
\textbf{Ingredient} & \textbf{Status} & \textbf{Derivation level} \\
\midrule
$\alpha = 1/137.036$ & Derived (zero parameters) & Theorem-grade \\
$A_f$ (9 prefactors) & Derived from $A_5$ operators & Constructive proof \\
$n_f$ (9 exponents) & Structurally assigned & Motivated, not theorem-grade \\
$v = 246.22$~GeV & Empirical input ($G_F$) & Measured (Formulation~I) \\
$v = M_P\alpha^8\sqrt{2\pi}$ & Derived, 0.05\% accurate & Tier 1.5 (Formulation~II) \\
\bottomrule
\end{tabular}
\end{center}

\noindent
We now expand on each of these assessments.

\paragraph{(a) The exponents $n_f$.}
The exponents $n_f = (k_f + k_H)/2$ are assigned from the spin$^c$
bundle degree of each fermion species on $\CP^2$.
The allowed values $\{0, 1, \tfrac{3}{2}, \tfrac{5}{2}\}$ correspond
to the available half-integer degrees, and the assignment to specific
fermions follows a pattern:
third-generation fermions carry the smallest exponent in their charge
sector, with each preceding generation incrementing by $\tfrac{1}{2}$
or $1$ depending on the sector.
This assignment is \emph{structurally motivated}---it follows a
consistent rule tied to the spin$^c$ bundle---but it has not been
derived as a theorem from an explicit construction of zero modes.
We regard the exponent sector as the weakest link in the derivation
chain and discuss its status explicitly in Section~\ref{sec:exponents}.

\paragraph{(b) The prefactors $A_f$.}
The nine prefactors are rational or algebraic numbers:
\begin{equation}\label{eq:prefactors}
  \{A_f\} = \left\{\tfrac{2}{3},\; 1,\; \sqrt{2},\;
  \tfrac{8}{3},\; 1,\; 1,\;
  6,\; \tfrac{6}{7},\; \tfrac{1}{42}\right\}.
\end{equation}
These are outputs of explicit operators ($G$, $K_d = J_3$,
$K_u = I_4$, $Q_d$, $D_\ell$) constructed from the $A_5$
(icosahedral) symmetry of the microsector.
The derivation is constructive: one writes down the operators,
evaluates them, and reads off the prefactors.
This part of the derivation is on solid ground.

\paragraph{(c) The Higgs VEV.}
Formulation~I uses $v = (\sqrt{2}\,G_F)^{-1/2} = 246.22$~GeV as
a single empirical input.
In this case the mass formula has one free parameter (or equivalently,
one mass is used to set the scale and the remaining eight are
predictions).
All mass \emph{ratios} remain parameter-free regardless.

Formulation~II uses the DFD-derived relation
$v = M_P\alpha^8\sqrt{2\pi} = 246.09$~GeV, which is accurate
to 0.05\% but whose derivation, while physically well-motivated
(the exponent $8 = \dim(\CP^2\times S^3) + 1$ is topologically
determined), has not been established at full mathematical rigor
comparable to the $\alpha$ derivation.

\paragraph{(d) What we do \emph{not} claim.}
\begin{itemize}
\item We do not claim theorem-grade derivations of the exponents
  $n_f$. The exponent assignments are consistent and structurally
  motivated, but an explicit zero-mode construction on $\CP^2$ that
  uniquely determines them remains an open problem.
\item We do not claim sub-ppm accuracy. The 1.4\% mean error
  reflects genuine approximations (quadrature discretization, leading-order
  spectral action, QCD running effects). This is qualitatively different
  from the $\alpha$ result, which is exact to all measured digits.
\item We do not claim that the $v = M_P\alpha^8\sqrt{2\pi}$ formula
  is proven at the same level as the $k_{\max} = 60$ result.
  We present it as a strong result (0.05\% accuracy from topology alone)
  whose rigorous proof from the microsector effective action remains
  to be written.
\end{itemize}

\subsection{What we \emph{do} claim}

Despite these caveats, the results are remarkable by any standard
in the mass-prediction literature:

\begin{enumerate}
\item \textbf{All nine mass ratios are parameter-free.}
  No Yukawa couplings are fitted. The ratio $m_c/m_t = \alpha$,
  the ratio $m_\tau/m_\mu = \sqrt{2}\,\alpha^{-1/2}$, and all
  other inter-fermion ratios follow from the formula with zero
  adjustable parameters.

\item \textbf{All nine absolute masses are reproduced to 1.4\%
  with at most one input.}
  Using only $G_F$ (equivalently $v$) as a single empirical input,
  or using $M_P$ with a derived $v$, the formula reproduces the
  entire charged fermion spectrum.

\item \textbf{The prefactors are algebraically determined.}
  The nine $A_f$ values are outputs of a definite operator algebra,
  not a fit. They are rational numbers (or $\sqrt{2}$) with
  specific group-theoretic origins.

\item \textbf{The mass hierarchy is topological.}
  The five-order-of-magnitude spread from $m_e$ to $m_t$ arises
  from powers of $\alpha \approx 1/137$, which is itself a
  topological invariant of $\CP^2$. The hierarchy is not fine-tuned.

\item \textbf{Every prediction is falsifiable.}
  With zero (or one) free parameters, improved measurements of
  any fermion mass can confirm or falsify the formula.
\end{enumerate}

\subsection{Comparison with prior work}

Previous approaches to fermion masses from geometry include:
\begin{itemize}
\item Texture zeros and discrete flavor symmetries, which reduce
  parameter counts but typically retain several free parameters;
\item Randall--Sundrum and extra-dimensional models, which generate
  hierarchies via wavefunction overlaps but require $O(10)$
  bulk mass parameters;
\item String landscape approaches, which predict mass distributions
  but not specific values.
\end{itemize}

The present approach differs in that the number of free parameters
is either one (Formulation~I) or zero (Formulation~II), and all
structure emerges from a fixed seven-dimensional manifold
$\CP^2\times S^3$.

\subsection{Outline}

Section~\ref{sec:formula} presents the master formula and its
ingredients.
Section~\ref{sec:exponents} discusses the exponent assignments
and their derivation status.
Section~\ref{sec:prefactors} derives the prefactors from
$A_5$ operator algebra.
Section~\ref{sec:vev} presents both formulations of the Higgs VEV.
Section~\ref{sec:results} gives the numerical results.
Section~\ref{sec:open} identifies the open problems.
Section~\ref{sec:conclusion} concludes.

%==========================================================================
\section{The Master Formula}
\label{sec:formula}
%==========================================================================

The charged fermion mass formula is:
\begin{equation}\label{eq:master2}
\boxed{m_f = A_f \times \alpha^{n_f} \times \frac{v}{\sqrt{2}}}
\end{equation}
with the following assignments:

\begin{table}[h]
\centering
\renewcommand{\arraystretch}{1.15}
\begin{tabular}{lcccrrc}
\toprule
Fermion & Sector & Gen & $A_f$ & $n_f$ & $m_{\rm pred}$ (MeV)
  & $m_{\rm PDG}$ (MeV) \\
\midrule
$e$   & lepton & 1 & $2/3$      & 2.5 & 0.528   & 0.511   \\
$\mu$ & lepton & 2 & 1          & 1.5 & 108.5   & 105.66  \\
$\tau$& lepton & 3 & $\sqrt{2}$ & 1.0 & 1796.7  & 1776.86 \\[3pt]
$u$   & up     & 1 & $8/3$      & 2.5 & 2.112   & 2.16    \\
$c$   & up     & 2 & 1          & 1.0 & 1270.5  & 1270    \\
$t$   & up     & 3 & 1          & 0   & 174\,100& 172\,760\\[3pt]
$d$   & down   & 1 & 6          & 2.5 & 4.752   & 4.67    \\
$s$   & down   & 2 & $6/7$      & 1.5 & 93.03   & 93.0    \\
$b$   & down   & 3 & $1/42$     & 0   & 4145    & 4180    \\
\bottomrule
\end{tabular}
\caption{The nine charged fermion masses from Eq.~\eqref{eq:master2}
  using $v/\!\sqrt{2} = 174.1$~GeV (Formulation~I).
  Mean $|$error$|$ = 1.42\%.}
\label{tab:masses}
\end{table}

%==========================================================================
\section{Exponent Assignments and Their Status}
\label{sec:exponents}
%==========================================================================

The exponents $n_f$ take values from the set
$\{0, 1, \tfrac{3}{2}, \tfrac{5}{2}\}$.
Within each charge sector, the assignment follows a
generation-descending rule:

\begin{center}
\begin{tabular}{lccc}
\toprule
& Gen~1 & Gen~2 & Gen~3 \\
\midrule
Leptons     & 5/2 & 3/2 & 1 \\
Up quarks   & 5/2 & 1   & 0 \\
Down quarks & 5/2 & 3/2 & 0 \\
\bottomrule
\end{tabular}
\end{center}

\paragraph{What is derived.}
The exponents originate from the spin$^c$ bundle on $\CP^2$:
$n_f = (k_f + k_H)/2$, where $k_f$ is the fermion's bundle degree
and $k_H$ is the Higgs bundle degree.
The available values are constrained by the topology---only
specific half-integer values are permitted.

\paragraph{What is not derived.}
The assignment of specific $k_f$ values to specific fermion species
(i.e., which fermion gets which bundle degree) has not been derived
from an explicit zero-mode construction.
Instead, it is assigned by requiring consistency with the observed
mass hierarchy:
first-generation fermions universally carry $n_f = 5/2$ (maximal
suppression), while third-generation fermions carry the smallest
exponent available in their sector.

\paragraph{Honest assessment.}
This is the least rigorous component of the derivation.
An explicit construction of fermion zero modes on $\CP^2$ that
uniquely determines the bundle-degree assignment would significantly
strengthen the result.
We note, however, that the number of exponent values is small (four),
the assignment pattern is systematic (not arbitrary), and the
resulting predictions are accurate, which suggests that the
assignment captures real structure.

%==========================================================================
\section{Prefactors from $A_5$ Operator Algebra}
\label{sec:prefactors}
%==========================================================================

The prefactors $A_f$ are the most solidly derived component of
the formula.
They emerge from explicit Yukawa operators:

\begin{itemize}
\item The \textbf{generation operator} $G = \mathrm{diag}(2/3, 1, 1)$,
  derived from the primed microsector trace:
  $G[1,1] = (\Tr\Pi - \Tr M_0)/\Tr\Pi = (9-3)/9 = 2/3$.

\item The \textbf{down-type kernel} $K_d = J_3$, the order-3
  permutation matrix from $A_5$.

\item The \textbf{up-type kernel} $K_u = I_4$, the order-4 element.

\item The \textbf{down-quark modifier} $Q_d = \mathrm{diag}(1, 6/7, 1/42)$.

\item The \textbf{lepton modifier} $D_\ell = \mathrm{diag}(1, 1, \sqrt{2})$.
\end{itemize}

The rational coefficients have specific origins:
the factor $6/7$ for the strange quark combines isospin and
generation-suppression effects;
the factor $1/42$ for the bottom quark is $1/(\alpha^{-1}/3)$
within the down-type operator.
The factor $\sqrt{2}$ for the tau lepton arises from the
third-generation lepton modifier.

The key point is that these are \emph{outputs} of a definite
algebraic construction, not inputs.
One does not tune the prefactors to match data; one constructs the
operators from $A_5$ representation theory and reads off the prefactors.

%==========================================================================
\section{The Higgs VEV: Two Formulations}
\label{sec:vev}
%==========================================================================

\subsection{Formulation I: $v$ from $G_F$ (one parameter)}

The Fermi constant $G_F = 1.1664 \times 10^{-5}$~GeV$^{-2}$ gives
$v = (\sqrt{2}\,G_F)^{-1/2} = 246.22$~GeV.
Using this measured value, the formula~\eqref{eq:master2} has
\textbf{one free parameter} (the overall scale $v$).
Equivalently, one fermion mass (e.g., $m_t$) sets the scale and the
remaining eight are predictions.

All mass \emph{ratios} are parameter-free in either formulation.

With this choice: mean $|$error$| = 1.42\%$, max error $= 3.32\%$ (electron).

\subsection{Formulation II: $v = M_P\alpha^8\sqrt{2\pi}$ (zero parameters)}

The DFD framework provides a derivation of the Higgs VEV:
\begin{equation}\label{eq:v_derived}
  v = M_P \times \alpha^8 \times \sqrt{2\pi} = 246.09~\text{GeV},
\end{equation}
where the exponent $8 = \dim(\CP^2\times S^3) + 1$ is topologically
determined.
This is accurate to 0.05\% versus the measured value.

Using $v_{\rm derived}$, the formula has \textbf{zero free parameters}:
the only input is $M_P$, and everything else---$\alpha$, $v$, all
$A_f$, all $n_f$---follows from the $\CP^2\times S^3$ topology.

With this choice: mean $|$error$| = 1.40\%$, max error $= 3.27\%$ (electron).

\paragraph{Caveat.}
The formula $v = M_P\alpha^8\sqrt{2\pi}$ is stated as Theorem~K.2
in the DFD corpus and attributed to microsector scaling.
The numerical agreement is striking, and the exponent is
topologically determined.
However, the complete proof from the microsector effective action
is less detailed than the proofs for $k_{\max} = 60$ or the
$A_f$ operators.
We classify it as ``Tier~1.5''---stronger than a motivated ansatz,
not yet a fully rigorous theorem.
Readers who prefer strict rigor should adopt Formulation~I (one
parameter); readers who accept the topological argument should
adopt Formulation~II (zero parameters).

%==========================================================================
\section{Numerical Results}
\label{sec:results}
%==========================================================================

\begin{table}[h]
\centering
\renewcommand{\arraystretch}{1.15}
\begin{tabular}{lrrrrr}
\toprule
& \multicolumn{2}{c}{Formulation I ($v_{\rm obs}$)}
& \multicolumn{2}{c}{Formulation II ($v_{\rm derived}$)} \\
\cmidrule(lr){2-3}\cmidrule(lr){4-5}
Fermion & $m_{\rm pred}$ & Error
  & $m_{\rm pred}$ & Error \\
\midrule
$e$   & 0.528 MeV  & $+3.32\%$ & 0.528 MeV  & $+3.27\%$ \\
$\mu$ & 108.5 MeV  & $+2.72\%$ & 108.5 MeV  & $+2.66\%$ \\
$\tau$& 1796.7 MeV & $+1.12\%$ & 1796.0 MeV & $+1.07\%$ \\
$u$   & 2.112 MeV  & $-2.23\%$ & 2.111 MeV  & $-2.27\%$ \\
$c$   & 1270.5 MeV & $+0.04\%$ & 1270.0 MeV & $-0.01\%$ \\
$t$   & 174.10 GeV & $+0.78\%$ & 174.01 GeV & $+0.73\%$ \\
$d$   & 4.752 MeV  & $+1.75\%$ & 4.750 MeV  & $+1.70\%$ \\
$s$   & 93.03 MeV  & $+0.03\%$ & 93.00 MeV  & $-0.02\%$ \\
$b$   & 4145 MeV   & $-0.83\%$ & 4143 MeV   & $-0.88\%$ \\
\midrule
\multicolumn{2}{l}{Mean $|$error$|$}
  & $\mathbf{1.42\%}$
  & & $\mathbf{1.40\%}$ \\
\multicolumn{2}{l}{Max $|$error$|$}
  & $3.32\%$
  & & $3.27\%$ \\
\bottomrule
\end{tabular}
\caption{Comparison of both formulations. Observed masses:
  PDG 2024 central values;
  quark masses at $\overline{\rm MS}$ running scales.}
\label{tab:comparison}
\end{table}

\paragraph{Highlights.}
\begin{itemize}
\item The charm and strange quarks are predicted to better than
  0.05\%---effectively exact within current uncertainties.
\item The top quark mass, $m_t = v/\!\sqrt{2} = 174.0$~GeV, follows
  from $A_t = 1$ and $n_t = 0$ with no further input.
  In Formulation~II this becomes
  $m_t = M_P\alpha^8\sqrt{\pi} = 174.0$~GeV.
\item The worst prediction is the electron ($+3.3\%$), which is
  the lightest fermion and most sensitive to the $\alpha^{5/2}$
  suppression factor.
\item No fermion deviates by more than 3.4\%.
\end{itemize}

%==========================================================================
\section{Open Problems}
\label{sec:open}
%==========================================================================

We identify three open problems whose resolution would bring the
mass derivation to the same level as the $\alpha$ derivation:

\begin{enumerate}
\item \textbf{Zero-mode construction for exponents.}
  An explicit construction of fermion zero modes on $\CP^2$ that
  uniquely determines which bundle degree $k_f$ each species carries.
  This would elevate the exponent assignments from ``structurally
  motivated'' to ``theorem-grade.''

\item \textbf{Rigorous proof of $v = M_P\alpha^8\sqrt{2\pi}$.}
  A complete derivation from the microsector effective action,
  at the level of rigor of the $k_{\max} = 60$ proof, establishing
  why the exponent is precisely $8 = \dim(\CP^2\times S^3) + 1$ and
  why the coefficient is precisely $\sqrt{2\pi}$.

\item \textbf{Understanding the residual errors.}
  The 1.4\% mean error likely arises from a combination of
  quadrature discretization, QCD running effects, and higher-order
  spectral action corrections.
  A systematic computation of these corrections could improve
  the accuracy and would test whether the formula is exact at some
  order or inherently approximate.
\end{enumerate}

%==========================================================================
\section{Conclusion}
\label{sec:conclusion}
%==========================================================================

We have presented a derivation of all nine charged fermion masses
from the geometry of $\CP^2\times S^3$, with the following
properties:

\begin{itemize}
\item All mass \textbf{ratios} are parameter-free predictions.
\item All absolute masses are reproduced to \textbf{1.4\% mean accuracy}
  with at most one empirical input ($G_F$), or with zero free
  parameters if the derived Higgs VEV is used.
\item The prefactors are \textbf{algebraically determined} from $A_5$
  operator algebra---not fitted.
\item The mass hierarchy spanning five orders of magnitude is
  \textbf{topological}: it arises from powers of $\alpha$, itself
  a topological invariant.
\end{itemize}

The derivation is not at the same level as the DFD
$\alpha$ formula, which achieves exact agreement with zero free
parameters and a complete proof chain.
The mass formula achieves 1.4\% accuracy, the exponent assignments
are structurally motivated rather than proven, and the
zero-parameter formulation relies on a VEV derivation that is
compelling but not yet fully rigorous.
These are clearly stated limitations, not hidden ones.

Nevertheless, by any standard in the fermion mass literature,
reproducing the entire charged fermion spectrum---including
the $m_t/m_e = 3.4\times 10^5$ hierarchy---from a single formula
with algebraically determined coefficients and at most one free
parameter represents significant progress toward a geometric
understanding of flavor.

\begin{thebibliography}{9}
\bibitem{DFD_alpha}
G.~T.~Alcock, ``Ab initio derivation of the fine-structure constant
from Density Field Dynamics,'' (2026).
\end{thebibliography}

\end{document}
